\documentclass[10pt,a4paper]{article}
\usepackage{amsmath,amssymb,amsthm,graphicx}
\title{NP-hardness of approximating sublinear-time conditional Kolmogorov complexity via secret sharing}
\author{Shuichi Hirahara}
\date{doi:10.4230/LIPIcs.CCC.2022.26; CC BY4.0}
\begin{document}
\maketitle
\noindent\textit{Editorial note.} Original author: Shuichi Hirahara. Official CCC2022 article26 (2022-07-11), DOI 10.4230/LIPIcs.CCC.2022.26. Creative Commons Attribution4.0.\par\medskip
\noindent\textit{Editorial note.} License: https://creativecommons.org/licenses/by/4.0/\par\medskip
\noindent\textit{Editorial note.} Faithful editable transcription of formal author Theorem6.6 and complete Section6.2 main proof, including the full generic secret-sharing reduction Theorem6.14, all Claims6.17-6.23, exact cited prior premises and every page25-29 continuation. The generic reduction and privacy support are uncounted context. Original mathematical expressions and CC BY4.0 license retained; no new proof or correction.\par\medskip
\noindent\textit{Allowed standard prerequisites (editorial).}\begin{itemize}
\item Efficient secret sharing for monotone access structures and MMSA gap hardness. Exact cited Benaloh-Leichter Lemma6.11 and Dinur-Safra/Dinur et al. Lemma6.13, including original author Proof Sketch, retained.
\item Amortized time-bounded Kolmogorov complexity approaches entropy. Exact prior Lemma6.15 citing [6,5] retained; author computational-depth corollary supplied in full.
\end{itemize}\medskip
\paragraph{3. Preliminaries: Notation.}
$[n]$ denotes $\{1,\ldots,n\}$. For a function $p:\mathbb N\to\mathbb N$ and $i\in\mathbb N$, the function $p^{(i)}:\mathbb N\to\mathbb N$ is recursively defined as follows: $p^{(0)}(n):=n$ and $p^{(i+1)}(n):=p^{(i)}(p(n))$ for every $n\in\mathbb N$.
\paragraph{Kolmogorov Complexity.}
For a string $x\in\{0,1\}^*$ and an index $i\in\mathbb N$, let $x_i$ denote the $i$-th bit of $x$ if $i\leq|x|$ and $\bot$ if $i>|x|$. We fix an efficient universal Turing machine $U$. Time-bounded Kolmogorov complexity is formally defined as follows.
\paragraph{Definition 3.1 (Time-bounded Kolmogorov complexity).}
For strings $x,y\in\{0,1\}^*$, a time bound $t\in\mathbb N\cup\{\infty\}$, and an oracle $A$, the $A$-oracle $t$-time-bounded Kolmogorov complexity of $x$ given $y$ is defined as
\[K^{t,A}(x\mid y):=\min\{|d|:U^{A,y,d}\text{ outputs }x_i\text{ on input }i\text{ in time }t\text{ for every }i\in[1,\ldots,|x|+1]\}.\]
Here, $U^{A,y,d}$ indicates that the universal Turing machine is given oracle access to $A$ and each bit of $y$ and $d\in\{0,1\}^*$. We omit the superscript $A$ if $A=\varnothing$, the superscript $t$ if $t=\infty$, and ``$\mid y$'' if $y$ is the empty string.
A simple counting argument implies the following basic fact of Kolmogorov-randomness.
\paragraph{Fact 3.2.}
For any integer $s\geq1$ and any string $y\in\{0,1\}^*$, the number of strings $x\in\{0,1\}^*$ such that $K(x\mid y)<s$ is less than $2^s$.
\begin{proof}
The number of programs of length less than $s$ is at most $\sum_{i=0}^{s-1}2^i<2^s$.
\end{proof}

\paragraph{Definition 3.5 (Time-Bounded Computational Depth [9,35]).}
For time bounds $s\in\mathbb N$ and $t\in\mathbb N\cup\{\infty\}$ and a string $x\in\{0,1\}^*$, the $(s,t)$-time-bounded computational depth of $x$ is defined as
\[\operatorname{cd}^{s,t}(x):=K^s(x)-K^t(x).\]
We omit the superscript $t$ if $t=\infty$.
\paragraph{Definition 6.1.}
For a polynomial $\tau:\mathbb N^3\to\mathbb N$, we define the promise problem $\mathsf{Gap}_{\tau}\mathsf{MINcKT}$ to be $(\Pi_{\mathrm{Yes}},\Pi_{\mathrm{No}})$ such that
\[\Pi_{\mathrm{Yes}}:=\{(x,y,1^t,1^s):K^t(x\mid y)\leq s-\operatorname{cd}^{t,\tau(|x|,|y|,t)}(y)\},\]
\[\Pi_{\mathrm{No}}:=\{(x,y,1^t,1^s):K^{\tau(|x|,|y|,t)}(x\mid y)>s+\log\tau(|x|,|y|,t)\}.\]
\paragraph{Definition 6.5.}
For a polynomial $\tau:\mathbb N^3\to\mathbb N$ and a function $\sigma:\mathbb N\to\mathbb N$, we define the promise problem $\mathsf{Gap}_{\tau,\sigma}\mathsf{MINcKT}$ to be $(\Pi_{\mathrm{Yes}},\Pi_{\mathrm{No}})$ such that
\[\Pi_{\mathrm{Yes}}:=\{(x,y,1^t,1^s):K^t(x\mid y)\leq s-\operatorname{cd}^{t,\tau(|x|,|y|,t)}(y)\},\]
\[\Pi_{\mathrm{No}}:=\{(x,y,1^t,1^s):K^{\tau(|x|,|y|,t)}(x\mid y)>\sigma(|x|)\cdot s+\log\tau(|x|,|y|,t)\}.\]
Note that $\mathsf{Gap}_{\tau,1}\mathsf{MINcKT}=\mathsf{Gap}_{\tau}\mathsf{MINcKT}$ is trivially reducible to $\mathsf{Gap}_{\tau,\sigma}\mathsf{MINcKT}$ for every $\sigma\geq1$.
\paragraph{Definition 6.7 (Access Structure).}
An access structure $\mathcal A\subseteq2^{[n]}$ is a ``monotone'' collection of subsets of $[n]$; that is, for every $T\supseteq S\in\mathcal A$, we have $T\in\mathcal A$. The minimum weight of $\mathcal A$ is defined to be $w(\mathcal A):=\min\{|T|:T\in\mathcal A\}$.
We will prove that there is a generic reduction from the problem of estimating the weight of access structures $\mathcal A$ to $\mathsf{GapMINcKT}$ if there exists an efficient secret sharing scheme for $\mathcal A$.
\paragraph{Definition 6.8 (Secret Sharing [12]).}
A secret sharing scheme for $\mathcal A$ is a pair $(\mathsf{Share},\mathsf{Rec})$ of a randomized algorithm $\mathsf{Share}$ and a deterministic algorithm $\mathsf{Rec}$ with the following properties for every $\ell\in\mathbb N$:
\begin{enumerate}
\item Correctness: For every $T\in\mathcal A$ and for every string $x\in\{0,1\}^\ell$, any output of $\mathsf{Share}(x)$ is a sequence $(y_1,\ldots,y_n)$ of $n$ strings that satisfies $\mathsf{Rec}(y_T)=x$, where $y_T:=\{(i,y_i):i\in T\}$.
\item Privacy: For every $T\notin\mathcal A$ and for every random variable $X$ on $\{0,1\}^\ell$, the random variables $X$ and $\mathsf{Share}(X)_T$ are statistically independent.
\end{enumerate}
We observe that the privacy condition can be stated in terms of Kolmogorov complexity. We mention in passing that Kolmogorov complexity-theoretic versions of privacy conditions are studied in, e.g., [10,48].
\paragraph{Definition 6.10.}
A family $\mathcal A=\{\mathcal A_\phi\}_{\phi\in\{0,1\}^*}$ of access structures is said to admit efficient secret sharing schemes if there exists a pair $(\mathsf{Share},\mathsf{Rec})$ of a randomized polynomial-time algorithm $\mathsf{Share}$ and a deterministic polynomial-time algorithm $\mathsf{Rec}$ such that for every $\phi\in\{0,1\}^*$, the pair $(\mathsf{Share}(\phi,-),\mathsf{Rec}(\phi,-))$ is a secret sharing scheme for the access structure $\mathcal A_\phi$.
Benaloh and Leichter [13] showed that access structures represented by monotone formulas admit efficient secret sharing schemes.
\paragraph{Lemma 6.11 ([13]).}
Let $\mathcal A:=\{\mathcal A_\phi\}_{\phi\in\{0,1\}^*}$ be the family of access structures $\mathcal A_\phi:=\{T\subseteq[n]:\phi(\chi_T)=1\}$, where $\phi$ is a monotone formula on $n$ variables and $\chi_T\in\{0,1\}^n$ denotes the characteristic vector of $T\subseteq[n]$. Then, $\mathcal A$ admits efficient secret sharing schemes.
\paragraph{Definition 6.12 (Minimum Monotone Satisfying Assignment; MMSA).}
For a monotone formula $\phi$ on $n$ variables, the weight of an assignment $\alpha\in\{0,1\}^n$ is defined to be $\sum_{i=1}^n\alpha_i$. Let $\mathsf{MMSA}(\phi)$ denote the minimum weight of $\alpha\in\{0,1\}^n$ such that $\phi(\alpha)=1$.
Observe that $\mathsf{MMSA}(\phi)=w(\mathcal A_\phi)$ for the family $\mathcal A$ of access structures of Lemma 6.11.

\paragraph{Theorem 6.6.}
Let $c>1$ be an arbitrary constant and $\tau:\mathbb N^3\to\mathbb N$ be a function such that $\tau(n,m,t)\leq n^c\cdot m^{1-1/c}\cdot t^c$ for all large $n,m$, and $t\in\mathbb N$. Then, $\mathsf{Gap}_{\tau,\sigma}\mathsf{MINcKT}$ is $\mathsf{NP}$-hard under one-query randomized polynomial-time reductions for some function $\sigma:\mathbb N\to\mathbb N$ such that $\sigma(n)=n^{1/(\log\log n)^{O(1)}}$.

\paragraph{Lemma 6.9.}
Let $(\mathsf{Share},\mathsf{Rec})$ be a secret sharing scheme for an access structure $\mathcal A$ over $[n]$. Then, for every $\ell$ and $k\in\mathbb N$, it holds that
\[\Pr\left[\min_{T\notin\mathcal A}K(X\mid\mathsf{Share}(X)_T)\geq\ell-n-k\right]\geq1-2^{-k},\]
where $X$ is the uniform distribution over $\{0,1\}^\ell$ and the probability is taken over $X$ as well as the internal randomness of $\mathsf{Share}$.
\begin{proof}
Fix an arbitrary subset $T\notin\mathcal A$ of $[n]$. Let $y_T\in\operatorname{supp}(\mathsf{Share}(X)_T)$ be an outcome of $\mathsf{Share}(X)_T$. By Fact 3.2, we obtain
\[\Pr[K(X\mid y_T)<\ell-n-k]\leq2^{-n-k}.\]
By the privacy of the secret sharing scheme, the random variables $X$ and $\mathsf{Share}(X)_T$ are statistically independent. By averaging the inequality above over all $y_T\in\operatorname{supp}(\mathsf{Share}(X)_T)$, we obtain
\[\Pr[K(X\mid\mathsf{Share}(X)_T)<\ell-n-k]\leq2^{-n-k}.\]
Finally, we take the union bound over all subsets $T\notin\mathcal A$ and conclude that
\[\Pr[\exists T\notin\mathcal A,K(X\mid\mathsf{Share}(X)_T)<\ell-n-k]\leq2^{-k}.\]
\end{proof}
It is known that $\mathsf{MMSA}$ is $\mathsf{NP}$-hard to approximate:
\paragraph{Lemma 6.13 ([22,21]).}
For some function $g(n)=n^{1/(\log\log n)^{O(1)}}$, it is $\mathsf{NP}$-hard to solve the promise problem $\mathsf{Gap}_g\mathsf{MMSA}=(\Pi_{\mathrm{Yes}},\Pi_{\mathrm{No}})$ defined as follows:
\[\Pi_{\mathrm{Yes}}:=\{(\phi,s):\mathsf{MMSA}(\phi)\leq s\},\qquad\Pi_{\mathrm{No}}:=\{(\phi,s):\mathsf{MMSA}(\phi)>s\cdot g(|\phi|)\},\]
where $|\phi|$ denotes the length of the binary string that represents $\phi$.
\begin{proof}[Proof Sketch]
Dinur and Safra [22] presented a polynomial-time reduction from any constraint satisfaction problem (CSP) to $\mathsf{Gap}_g\mathsf{MMSA}$ with the following property: Let $\psi$ be an instance of the CSP. Let $(\phi,s)$ be the output of the reduction. Let $D$ be the arity of $\psi$. If $\psi$ is satisfiable, then $\mathsf{MMSA}(\phi)\leq s$. If there is no assignment satisfying a $1/(2(2Dg)^D)$ fraction of the constraints of $\psi$, then $\mathsf{MMSA}(\phi)>s\cdot g$, where $g$ is an arbitrary parameter. Dinur et al. [21] showed that it is $\mathsf{NP}$-hard to decide whether a given CSP with $n$ variables and of arity $D=(\log\log n)^{O(1)}$ is satisfiable or every assignment satisfies at most a $1/\operatorname{poly}(n)$ fraction of the constraints. Combining these two hardness results, we conclude that $\mathsf{Gap}_g\mathsf{MMSA}$ is $\mathsf{NP}$-hard, where $g\geq\operatorname{poly}(n)^{1/D}\cdot D^{-D}\geq n^{1/(\log\log n)^{O(1)}}$.
\end{proof}
\paragraph{Lemma 6.15 ([6,5]).}
Let $D=\{D_x\}_{x\in\{0,1\}^*}$ be a family of distributions sampled by circuits of size $\operatorname{poly}(|x|)$. Then, there exist a polynomial $p$ and a constant $\delta>0$ such that for every $x\in\{0,1\}^*$ and for every $D\geq p(|x|)$,
\[\frac1D\cdot K^t(y^1,\ldots,y^D)\leq H(D_x)+\frac12\cdot D^{-\delta},\]
\[\frac1D\cdot K(y^1,\ldots,y^D)\geq H(D_x)-\frac12\cdot D^{-\delta}\]
holds with probability at least $1-2^{-|x|}$ over a random choice of $D$ independent samples $y^1,\ldots,y^D$ from $D_x$, where $t:=p(|x|)$.
This lemma immediately implies that the computational depth of independent samples from $D_x$ is small.
\paragraph{Corollary 6.16.}
Under the same assumptions as Lemma 6.15, with probability at least $1-2^{-|x|}$ over a random choice of $D$ independent samples $y^1,\ldots,y^D$ from $D_x$, it holds that
\[\operatorname{cd}^t(y^1,\ldots,y^D)\leq D^{1-\delta}.\]
\begin{proof}
It follows from Lemma 6.15 that
\[\operatorname{cd}^t(y^1,\ldots,y^D)=K^t(y^1,\ldots,y^D)-K(y^1,\ldots,y^D)\leq D^{1-\delta}.\]
\end{proof}

\paragraph{6.2.3. A Generic Connection from Secret Sharing Schemes.}
We now formally state a generic reduction that reduces the problem of estimating the weight of access structures to $\mathsf{GapMINcKT}$.
\paragraph{Theorem 6.14.}
Let $\mathcal A=\{\mathcal A_\phi\}_{\phi\in\{0,1\}^*}$ be a family of access structures that admits efficient secret sharing schemes. Let $\tau:\mathbb N^3\to\mathbb N$ be a function such that $t\leq\tau(n,m,t)\leq n^c\cdot m^{1-1/c}\cdot t^c$ for some constant $c>1$. Then, there exists a randomized polynomial-time algorithm $R$ that takes $\phi\in\{0,1\}^*$ as input and outputs $x,y\in\{0,1\}^*$ and $t,\rho\in\mathbb N$ such that
\[\frac\rho{8c}\cdot w(\mathcal A_\phi)+\log t'\leq K^{t'}(x\mid y)\leq K^t(x\mid y)\leq\rho\cdot w(\mathcal A_\phi)-\operatorname{cd}^t(y)\]
holds with probability at least $1-o(1)$ over the internal randomness of $R$, where $t':=\tau(|x|,|y|,t)$.
It is easy to observe that this reduction implies $\mathsf{NP}$-hardness of $\mathsf{GapMINcKT}$.
\begin{proof}[Proof of Theorem 6.6]
Let $\mathcal A=\{\mathcal A_\phi\}_\phi$ be the family of access structures for monotone formulas given in Lemma 6.11. Let $R$ be the reduction of Theorem 6.14. Consider a reduction $R'$ that reduces an instance $(\phi,s)$ of $\mathsf{Gap}_g\mathsf{MMSA}$ to an instance $(x,y,1^t,1^{s'})$ of $\mathsf{Gap}_{\tau}\mathsf{MINcKT}$ such that $(x,y,t,\rho):=R(\phi)$ and $s':=s\cdot\rho$.
Let $(\Pi_{\mathrm{Yes}},\Pi_{\mathrm{No}}):=\mathsf{Gap}_g\mathsf{MMSA}$. We claim below that the reduction $R'$ reduces $\mathsf{Gap}_g\mathsf{MMSA}$ to $\mathsf{Gap}_{\tau,\sigma}\mathsf{MINcKT}$ for some $\sigma$, which implies $\mathsf{NP}$-hardness of $\mathsf{Gap}_{\tau,\sigma}\mathsf{MINcKT}$ by Lemma 6.13. If $(\phi,s)\in\Pi_{\mathrm{Yes}}$, then $w(\mathcal A_\phi)=\mathsf{MMSA}(\phi)\leq s$. By the properties of $R$, with high probability, it holds that
\[K^t(x\mid y)\leq\rho\cdot w(\mathcal A_\phi)-\operatorname{cd}^t(y)\leq s'-\operatorname{cd}^{t,t'}(y),\]
which implies that $(x,y,1^t,1^{s'})$ is a Yes instance of $\mathsf{Gap}_{\tau,\sigma}\mathsf{MINcKT}$. If $(\phi,s)\in\Pi_{\mathrm{No}}$, then $w(\mathcal A_\phi)=\mathsf{MMSA}(\phi)>g(|\phi|)\cdot s$. Also by the properties of $R$, with high probability, it holds that
\[K^{t'}(x\mid y)\geq\frac\rho{8c}\cdot w(\mathcal A_\phi)+\log t'>\frac{g(|\phi|)}{8c}\cdot s'+\log t',\]
which implies that $(x,y,1^t,1^{s'})$ is a No instance of $\mathsf{Gap}_{\tau,\sigma}\mathsf{MINcKT}$ for some $\sigma(|x|):=g(|\phi|)/(8c)\geq|\phi|^{1/(\log\log|\phi|)^{O(1)}}\geq|x|^{1/(\log\log|x|)^{O(1)}}$, where the last inequality follows from the fact that $|\phi|\geq|x|^{\Omega(1)}$.
\end{proof}
It remains to prove Theorem 6.14. To show that the additive error $\operatorname{cd}^t(y)$ is relatively small, we use the fact that for $D$ independent random samples $y^1,\ldots,y^D$ from a distribution $D$ samplable by polynomial-size circuits, the amortized time-bounded Kolmogorov complexity of $y^1,\ldots,y^D$ approaches the entropy of $D$ asymptotically.
We are now ready to prove Theorem 6.14.
\begin{proof}[Proof of Theorem 6.14]
Fix an input $\phi\in\{0,1\}^*$. Let $n$ be the number of parties in the access structure $\mathcal A_\phi$. Let $\lambda=O(\log|\phi|)$, $t$, and $\ell$ be parameters chosen later.
We first define a family $D=\{D_\phi\}_{\phi\in\{0,1\}^*}$ of distributions by using the following sampling procedure: Choose $x\sim\{0,1\}^\ell$ uniformly at random. Let $(s_1,\ldots,s_n):=\mathsf{Share}(\phi,x)$ and let $m$ be the length of each share; i.e., $m:=|s_i|$. Pick $k_1,\ldots,k_n\sim\{0,1\}^\lambda$ randomly. We define a string $y\in\{0,1\}^{2^\lambda\cdot m}$, which we identify with a function $y:\{0,1\}^\lambda\to\{0,1\}^m$: For every $q\in\{0,1\}^\lambda$, let $y(q):=s_i$ if $q=k_i$ for some $i\in[n]$ and let $y(q):=0^m$ otherwise. Note that $y$ is not well defined if $(k_i)_{i\in[n]}$ is not pairwise distinct; however, we will show that $y$ is well defined with high probability. The output of the sampling procedure is defined as $(x,y,k,s)$, where $k:=(k_1,\ldots,k_n)$ and $s:=(s_1,\ldots,s_n)$.
The algorithm $R$ operates as follows: For a given input $\phi\in\{0,1\}^*$, pick $D$ independent samples $(x^1,y^1,k^1,s^1),\ldots,(x^D,y^D,k^D,s^D)$ from $D_\phi$. Let $x:=(x^1,\ldots,x^D)$, $y:=(y^1,\ldots,y^D)$, and $\rho:=2\lambda D$. The output of $R$ is defined to be $(x,y,t,\rho)$.
Below, we prove the correctness of $R$ using a sequence of claims. Let $k:=(k_i^d:i\in[n],d\in[D])$ and let $s:=(s_i^d:i\in[n],d\in[D])$. We first observe that $k$ is Kolmogorov-random with high probability:
\paragraph{Claim 6.17.}
With probability at least $1-o(1)$, it holds that
\[K(k\mid s)\geq nD\lambda-\log n.\tag{16}\]
\begin{proof}
Since $k=(k_i^d)_{i\in[n],d\in[D]}$ is uniformly distributed over $(\{0,1\}^\lambda)^{nD}$ and independent of $s$, the claim follows from Fact 3.2.
\end{proof}
\paragraph{Claim 6.18.}
If Equation (16) holds, then $k=(k_i^d)_{i\in[n],d\in[D]}$ is pairwise distinct and hence $y$ is well defined.
\begin{proof}
If there exists $(i,d)\neq(i',d')$ such that $k_i^d=k_{i'}^{d'}$, then $k$ can be described by a program of size $nD\lambda-\lambda+O(\log nD)$. Choosing a large enough $\lambda=O(\log|\phi|)$, this implies that $K(k)\leq nD\lambda-\lambda+O(\log nD)<nD\lambda-\log n$, which contradicts Equation (16).
\end{proof}
Below, we assume that Equation (16) holds. We prove an upper bound of $K^t(x\mid y)$.
\paragraph{Claim 6.19.}
$K^t(x\mid y)\leq\rho\cdot w(\mathcal A_\phi)-\operatorname{cd}^t(y)$ holds with probability at least $1-o(1)$.
\begin{proof}
Let $T\in\mathcal A_\phi$ be a minimum authorized set of parties such that $|T|=w(\mathcal A_\phi)$. We present an oracle program $M^y$ that takes an index $d$ and outputs $x^d\in\{0,1\}^\ell$ as follows: $M^y$ takes $d\in[D]$ as input and $T$, $\{k_i^d:i\in T,d\in[D]\}$ and $\phi$ as hard-wired input, computes $\{s_i^d:i\in T\}=\{y^d(k_i^d):i\in T\}$ by making queries to the oracle $y$, and outputs $\mathsf{Rec}(\phi,s_T^d)$, which is equal to $x^d$ by the correctness of a secret sharing scheme and the assumption that $T\in\mathcal A_\phi$. The size of the oracle machine $M$ is at most
\[O(|T|\cdot\log n)+|T|\cdot D\cdot\lambda+|\phi|+O(\log D)\leq2\lambda D\cdot w(\mathcal A_\phi)-D^{1-\delta},\]
where $\delta>0$ is the constant from Corollary 6.16 and this inequality follows by choosing sufficiently large $D\geq\operatorname{poly}(|\phi|)$. Applying Corollary 6.16 to the distribution of $y$, we obtain $\operatorname{cd}^t(y)\leq D^{1-\delta}$ with probability at least $1-o(1)$; thus, $|M|\leq\rho\cdot w(\mathcal A_\phi)-\operatorname{cd}^t(y)$. The running time of $M$ is at most $\operatorname{poly}(|\phi|)$, which is independent of $\lambda$. We choose $t$ so that this running time is at most $t$.
\end{proof}
The remainder of the proof is devoted to proving the lower bound.
\paragraph{Claim 6.20.}
Let $\theta:=\frac\rho{8c}\cdot w(\mathcal A_\phi)+\log t'$. Then, $K^{t'}(x\mid y)\geq\theta$ with probability at least $1-o(1)$.
We first clarify the condition of random variables under which the claim holds. Let $s_{[n]}^{[D]\setminus\{d\}}:=(s_i^{d'}:i\in[n],d'\in[D]\setminus\{d\})$.
\paragraph{Claim 6.21.}
With probability at least $1-o(1)$, it holds that for every unauthorized set $T\notin\mathcal A_\phi$ of parties and for every $d\in[D]$,
\[K(x^d\mid s_T^d,s_{[n]}^{[D]\setminus\{d\}},k)\geq\ell-n-2\log D.\tag{17}\]
\begin{proof}
Fix $d\in[D]$. Since the random variable $x^d$ is independent of $s_{[n]}^{[D]\setminus\{d\}}$ and $k$, by Lemma 6.9,
\[\min_{T\notin\mathcal A_\phi}K(x^d\mid s_T^d,s_{[n]}^{[D]\setminus\{d\}},k)\geq\ell-n-2\log D\]
holds with probability at least $1-1/D^2$. The claim follows by taking a union bound over all $d\in[D]$.
\end{proof}
In what follows, we assume Equations (16) and (17) and prove $K^{t'}(x\mid y)\geq\theta$, which will complete the proof of Claim 6.20. Assume, by way of contradiction, that $K^{t'}(x\mid y)<\theta$. Let $M\in\{0,1\}^*$ be the description of an oracle program such that $|M|\leq\theta+O(\log\ell)$ and $M^y(d)$ outputs $x^d\in\{0,1\}^\ell$ in time $t'':=t'\ell$ for every $d\in[D]$; the machine $M$ can be constructed from a program that witnesses $K^{t'}(x\mid y)<\theta$. For each $d\in[D]$, let $T(M,d)$ be the set of indices $i\in[n]$ such that some bit of $y^d(k_i^d)\in\{0,1\}^m$ is queried during the computation of $M^y(d)$.
\paragraph{Claim 6.22.}
Let $M$ be an oracle program that runs in time $t''$ and let $\alpha:=\sum_{d=1}^D|T(M,d)|$. Then,
\[K(k\mid s)\leq|M|+(nD-\alpha)\cdot\lambda+\alpha\cdot(\log t''+\log nD)+O(\log nD).\]
\begin{proof}
By the definition of $T(M,d)$, for every $i\in T(M,d)$, there exists a time step $t_{i,d}\in[t'']$ such that $M^y$ makes a query $k_i^d$ to the oracle $y$ on input $d\in[D]$. To describe $k$ given $s$, consider the following program $M'$: $M'$ takes
\[\{(i,d,t_{i,d})\in[n]\times[D]\times[t'']:i\in T(M,d)\},\]
\[\{k_i^d\in\{0,1\}^\lambda:(i,d)\in[n]\times[D],\ i\notin T(M,d)\},\]
and $M$ as input and simulates $M^y$ on input $d$ for every $d\in[D]$. At the time step $t_{i,d}$ of the simulation of $M^y(d)$ for some $i\in T(M,d)$, $M'$ reads the query $k_i^d$ that $M^y(d)$ makes to the oracle $y^d$ and answers the query with $s_i^d$. $M'$ continues the simulation until the time step $t''$, at which point $M'$ knows $k_i^d$ for every $i\in T(M,d)$. Finally, $M'$ outputs $k$. The input of $M'$ can be encoded as a string of length $\alpha\cdot(\log t''+\log nD)+(nD-\alpha)\cdot\lambda+|M|+O(\log nD)$.
\end{proof}
It follows from Claim 6.22 and Equation (16) that
\[nD\lambda-\log n\leq|M|+(nD-\alpha)\cdot\lambda+\alpha\cdot(\log t''+\log nD)+O(\log nD),\]
which can be simplified to
\[(\lambda-\log t''-\log nD)\cdot\alpha\leq|M|+O(\log nD).\tag{18}\]
Since $t''/\ell=t'=\tau(|x|,|y|,t)\leq|x|^c\cdot|y|^{1-1/c}\cdot t^c\leq2^{(1-1/c)\cdot\lambda}\cdot(D\cdot|\phi|\cdot\ell\cdot m)^{O(c)}$, we may choose $\lambda=O(\log|\phi|)$ large enough so that $t''\leq2^{(1-1/2c)\cdot\lambda-\log nD}$. Then, by Equation (18), we obtain
\begin{align*}
\alpha&\leq\frac{2c}{\lambda}\cdot(|M|+O(\log nD))\\
&<\frac{4c}{\lambda}\cdot(\theta-\log t')=\frac\rho{2\lambda}\cdot w(\mathcal A_\phi)=D\cdot w(\mathcal A_\phi),
\end{align*}
where the second inequality follows from $|M|+O(\log nD)\leq\theta+O(\log nD\ell t')<2\theta$.\footnote{We may assume without loss of generality that $w(\mathcal A_\phi)\geq1$; then, we have $\theta\geq\Omega(\rho)\geq\Omega(D)$, which can be assumed to be larger than logarithmic terms.} It follows that there exists $d\in[D]$ such that $|T(M,d)|\leq\alpha/D<w(\mathcal A_\phi)$, which implies that
\[T(M,d)\notin\mathcal A_\phi.\]
By Equation (17), we obtain
\[K(x^d\mid s_{T(M,d)}^d,s_{[n]}^{[D]\setminus\{d\}},k)\geq\ell-n-2\log D.\]
However, this contradicts the following claim.
\paragraph{Claim 6.23.}
Let $M$ be an oracle program such that $M^y(d)$ outputs $x^d$ for every $d\in[D]$. Then, for every $d\in[D]$, it holds that
\[K(x^d\mid s_{T(M,d)}^d,s_{[n]}^{[D]\setminus\{d\}},k)\leq|M|+O(\log D).\]
\begin{proof}
Since $M^y(d)$ outputs $x^d$ without making any query in $\{k_i^d:i\in[n]\setminus T(M,d)\}$, we can define another oracle $z=(z^1,\ldots,z^D)$ such that $z^{d'}:=y^{d'}$ for every $d'\in[D]\setminus\{d\}$ and $z^d(k_i^d):=s_i^d$ for every $i\in T(M,d)$ and $z^d(k):=0^m$ otherwise, so that there is no difference between $y$ and $z$ on inputs queried by $M^y(d)$; therefore, we obtain $x^d=M^y(d)=M^z(d)$. The claim follows by observing that the oracle $z$ can be constructed from $M$, $d\in[D]$, $s_{T(M,d)}^d$, $s_{[n]}^{[D]\setminus\{d\}}$, and $k$.
\end{proof}
We conclude that $\ell-n-2\log D\leq|M|+O(\log D)\leq O(\theta)\leq O(n\lambda D)$, which is a contradiction by letting $\ell\gg O(n\lambda D)$. This completes the proof of Claim 6.20.
\end{proof}
\paragraph{Remark 6.24.}
It is not hard to observe that our NP-hardness hardness reduction also proves the NP-hardness of $\mathsf{McKTP}$ [4] because the upper bound of Claim 6.19 also holds for $KT(x\mid y)$.

\begin{thebibliography}{99}
\bibitem[1]{published-1} Dorit Aharonov and Oded Regev. Lattice problems in NP $\cap$ coNP. J. ACM, 52(5):749–765, 2005. doi:10.1145/1089023.1089025.
\bibitem[2]{published-2} Miklós Ajtai. Generating Hard Instances of Lattice Problems (Extended Abstract). In Proceedings of the Symposium on the Theory of Computing (STOC), pages 99–108, 1996. doi:10.1145/237814.237838.
\bibitem[3]{published-3} Eric Allender, Harry Buhrman, Michal Koucký, Dieter van Melkebeek, and Detlef Ronneburger. Power from Random Strings. SIAM J. Comput., 35(6):1467–1493, 2006. doi:10.1137/ 050628994.
\bibitem[4]{published-4} Eric Allender, Mahdi Cheraghchi, Dimitrios Myrisiotis, Harsha Tirumala, and Ilya Volkovich. One-Way Functions and a Conditional Variant of MKTP. In Proceedings of the Foundations of Software Technology and Theoretical Computer Science (FSTTCS), pages 7:1–7:19, 2021. doi:10.4230/LIPIcs.FSTTCS.2021.7.
\bibitem[5]{published-5} Eric Allender, John Gouwar, Shuichi Hirahara, and Caleb Robelle. Cryptographic Hard- ness Under Projections for Time-Bounded Kolmogorov Complexity. In Proceedings of the International Symposium on Algorithms and Computation (ISAAC), pages 54:1–54:17, 2021. doi:10.4230/LIPIcs.ISAAC.2021.54.
\bibitem[6]{published-6} Eric Allender, Joshua A. Grochow, Dieter van Melkebeek, Cristopher Moore, and Andrew Morgan. Minimum Circuit Size, Graph Isomorphism, and Related Problems. SIAM J. Comput., 47(4):1339–1372, 2018. doi:10.1137/17M1157970.
\bibitem[7]{published-7} Eric Allender, Lisa Hellerstein, Paul McCabe, Toniann Pitassi, and Michael E. Saks. Minimizing Disjunctive Normal Form Formulas and AC0 Circuits Given a Truth Table. SIAM J. Comput., 38(1):63–84, 2008. doi:10.1137/060664537.
\bibitem[8]{published-8} Eric Allender, Michal Koucký, Detlef Ronneburger, and Sambuddha Roy. The pervasive reach of resource-bounded Kolmogorov complexity in computational complexity theory. J. Comput. Syst. Sci., 77(1):14–40, 2011. doi:10.1016/j.jcss.2010.06.004.
\bibitem[9]{published-9} Luis Antunes, Lance Fortnow, Dieter van Melkebeek, and N. V. Vinodchandran. Computational depth: Concept and applications. Theor. Comput. Sci., 354(3):391–404, 2006. doi:10.1016/ j.tcs.2005.11.033.
\bibitem[10]{published-10} Luis Antunes, Sophie Laplante, Alexandre Pinto, and Liliana C. M. Salvador. Cryptographic Security of Individual Instances. In Proceedings of Information Theoretic Security (ICITS), pages 195–210, 2007. doi:10.1007/978-3-642-10230-1\_17.
\bibitem[11]{published-11} Luis Filipe Coelho Antunes and Lance Fortnow. Worst-Case Running Times for Average-Case Algorithms. In Proceedings of the Conference on Computational Complexity (CCC), pages 298–303, 2009. doi:10.1109/CCC.2009.12.
\bibitem[12]{published-12} Amos Beimel. Secret-Sharing Schemes: A Survey. In The Third International Workshop on Coding and Cryptology (IWCC), pages 11–46, 2011. doi:10.1007/978-3-642-20901-7\_2.
\bibitem[13]{published-13} Josh Cohen Benaloh and Jerry Leichter. Generalized Secret Sharing and Monotone Functions. In Proceedings of the International Cryptology Conference (CRYPTO), pages 27–35, 1988. doi:10.1007/0-387-34799-2\_3.
\bibitem[14]{published-14} George Robert Blakley. Safeguarding cryptographic keys. In International Workshop on Managing Requirements Knowledge (MARK), pages 313–318, 1979. doi:10.1109/MARK.1979. 8817296.
\bibitem[15]{published-15} Manuel Blum and Sampath Kannan. Designing Programs that Check Their Work. J. ACM, 42(1):269–291, 1995. doi:10.1145/200836.200880.
\bibitem[16]{published-16} Andrej Bogdanov and Luca Trevisan. Average-Case Complexity. Foundations and Trends in Theoretical Computer Science, 2(1), 2006. doi:10.1561/0400000004.
\bibitem[17]{published-17} Andrej Bogdanov and Luca Trevisan. On Worst-Case to Average-Case Reductions for NP Problems. SIAM J. Comput., 36(4):1119–1159, 2006. doi:10.1137/S0097539705446974.
\bibitem[18]{published-18} Harry Buhrman, Lance Fortnow, and Aduri Pavan. Some Results on Derandomization. Theory Comput. Syst., 38(2):211–227, 2005. doi:10.1007/s00224-004-1194-y.
\bibitem[19]{published-19} Harry Buhrman and Elvira Mayordomo. An Excursion to the Kolmogorov Random Strings. J. Comput. Syst. Sci., 54(3):393–399, 1997. doi:10.1006/jcss.1997.1484.
\bibitem[20]{published-20} Lijie Chen, Shuichi Hirahara, and Neekon Vafa. Average-case Hardness of NP and PH from Worst-case Fine-grained Assumptions. In Proceedings of the Innovations in Theoretical Computer Science Conference (ITCS), pages 67:1–67:17, 2022.
\bibitem[21]{published-21} Irit Dinur, Prahladh Harsha, and Guy Kindler. Polynomially Low Error PCPs with polyloglog n Queries via Modular Composition. In Proceedings of the Symposium on Theory of Computing (STOC), pages 267–276, 2015. doi:10.1145/2746539.2746630.
\bibitem[22]{published-22} Irit Dinur and Shmuel Safra. On the hardness of approximating label-cover. Inf. Process. Lett., 89(5):247–254, 2004. doi:10.1016/j.ipl.2003.11.007.
\bibitem[23]{published-23} Joan Feigenbaum and Lance Fortnow. Random-Self-Reducibility of Complete Sets. SIAM J. Comput., 22(5):994–1005, 1993. doi:10.1137/0222061.
\bibitem[24]{published-24} Lance Fortnow and Martin Kummer. On Resource-Bounded Instance Complexity. Theor. Comput. Sci., 161(1\&2):123–140, 1996. doi:10.1016/0304-3975(95)00097-6.
\bibitem[25]{published-25} Lance Fortnow, Troy Lee, and Nikolai K. Vereshchagin. Kolmogorov Complexity with Error. In STACS 2006, 23rd Annual Symposium on Theoretical Aspects of Computer Science, Marseille, France, February 23-25, 2006, Proceedings, pages 137–148, 2006. doi:10.1007/11672142\_10.
\bibitem[26]{published-26} Lance Fortnow and Nick Reingold. PP is Closed Under Truth-Table Reductions. Inf. Comput., 124(1):1–6, 1996. doi:10.1006/inco.1996.0001.
\bibitem[27]{published-27} Eran Gat and Shafi Goldwasser. Probabilistic Search Algorithms with Unique Answers and Their Cryptographic Applications. Electron. Colloquium Comput. Complex., page 136, 2011. URL: https://eccc.weizmann.ac.il/report/2011/136.
\bibitem[28]{published-28} Halley Goldberg and Valentine Kabanets. A Simpler Proof of the Worst-Case to Average-Case Reduction for Polynomial Hierarchy via Symmetry of Information. Electronic Colloquium on Computational Complexity (ECCC), 007, 2022.
\bibitem[29]{published-29} Oded Goldreich and Shafi Goldwasser. On the Limits of Nonapproximability of Lattice Problems. J. Comput. Syst. Sci., 60(3):540–563, 2000. doi:10.1006/jcss.1999.1686.
\bibitem[30]{published-30} Ishay Haviv and Oded Regev. Tensor-based Hardness of the Shortest Vector Problem to within Almost Polynomial Factors. Theory of Computing, 8(1):513–531, 2012. doi:10.4086/ toc.2012.v008a023.
\bibitem[31]{published-31} Shuichi Hirahara. Non-black-box Worst-case to Average-case Reductions within NP. In Proceedings of the Symposium on Foundations of Computer Science (FOCS), pages 247–258, 2018.
\bibitem[32]{published-32} Shuichi Hirahara. Characterizing Average-Case Complexity of PH by Worst-Case Meta- Complexity. In Proceedings of the Symposium on Foundations of Computer Science (FOCS), pages 50–60, 2020.
\bibitem[33]{published-33} Shuichi Hirahara. Non-Disjoint Promise Problems from Meta-Computational View of Pseudor- andom Generator Constructions. In Proceedings of the Computational Complexity Conference (CCC), pages 20:1–20:47, 2020. doi:10.4230/LIPIcs.CCC.2020.20.
\bibitem[34]{published-34} Shuichi Hirahara. Unexpected hardness results for Kolmogorov complexity under uniform reductions. In Proceedings of the Symposium on Theory of Computing (STOC), pages 1038– 1051, 2020. doi:10.1145/3357713.3384251.
\bibitem[35]{published-35} Shuichi Hirahara. Average-case hardness of NP from exponential worst-case hardness assump- tions. In Proceedings of the Symposium on Theory of Computing (STOC), pages 292–302, 2021. doi:10.1145/3406325.3451065.
\bibitem[36]{published-36} Shuichi Hirahara and Mikito Nanashima. On Worst-Case Learning in Relativized Heuristica. In Proceedings of the Symposium on Foundations of Computer Science (FOCS), 2021.
\bibitem[37]{published-37} Shuichi Hirahara, Igor Carboni Oliveira, and Rahul Santhanam. NP-hardness of Minimum Circuit Size Problem for OR-AND-MOD Circuits. In Proceedings of the Computational Complexity Conference (CCC), pages 5:1–5:31, 2018. doi:10.4230/LIPIcs.CCC.2018.5.
\bibitem[38]{published-38} Shuichi Hirahara and Rahul Santhanam. Errorless versus Error-prone Average-case Complexity. In Proceedings of the Innovations in Theoretical Computer Science Conference (ITCS), pages 38:1–38:23, 2022.
\bibitem[39]{published-39} Shuichi Hirahara and Osamu Watanabe. Limits of Minimum Circuit Size Problem as Oracle. In Proceedings of the Conference on Computational Complexity (CCC), pages 18:1–18:20, 2016. doi:10.4230/LIPIcs.CCC.2016.18.
\bibitem[40]{published-40} Rahul Ilango. Approaching MCSP from Above and Below: Hardness for a Conditional Variant and AC0 [p]. In Proceedings of the Innovations in Theoretical Computer Science Conference (ITCS), pages 34:1–34:26, 2020. doi:10.4230/LIPIcs.ITCS.2020.34.
\bibitem[41]{published-41} Rahul Ilango. Constant Depth Formula and Partial Function Versions of MCSP are Hard. In Proceedings of the Symposium on Foundations of Computer Science (FOCS), pages 424–433, 2020.
\bibitem[42]{published-42} Rahul Ilango, Bruno Loff, and Igor Carboni Oliveira. NP-Hardness of Circuit Minimization for Multi-Output Functions. In Proceedings of the Computational Complexity Conference (CCC), pages 22:1–22:36, 2020. doi:10.4230/LIPIcs.CCC.2020.22.
\bibitem[43]{published-43} Russell Impagliazzo. A Personal View of Average-Case Complexity. In Proceedings of the Structure in Complexity Theory Conference, pages 134–147, 1995. doi:10.1109/SCT.1995. 514853.
\bibitem[44]{published-44} Russell Impagliazzo. Relativized Separations of Worst-Case and Average-Case Complexities for NP. In Proceedings of the Conference on Computational Complexity (CCC), pages 104–114, 2011. doi:10.1109/CCC.2011.34.
\bibitem[45]{published-45} Russell Impagliazzo and Leonid A. Levin. No Better Ways to Generate Hard NP Instances than Picking Uniformly at Random. In Proceedings of the Symposium on Foundations of Computer Science (FOCS), pages 812–821, 1990. doi:10.1109/FSCS.1990.89604.
\bibitem[46]{published-46} Russell Impagliazzo and Avi Wigderson. P = BPP if E Requires Exponential Circuits: Derandomizing the XOR Lemma. In Proceedings of the Symposium on the Theory of Computing (STOC), pages 220–229, 1997. doi:10.1145/258533.258590.
\bibitem[47]{published-47} Valentine Kabanets and Jin-yi Cai. Circuit minimization problem. In Proceedings of the Symposium on Theory of Computing (STOC), pages 73–79, 2000. doi:10.1145/335305. 335314.
\bibitem[48]{published-48} Tarik Kaced. Almost-perfect secret sharing. In Proceedings of the International Symposium on Information Theory (ISIT), pages 1603–1607, 2011. doi:10.1109/ISIT.2011.6033816.
\bibitem[49]{published-49} Mauricio Karchmer and Avi Wigderson. On Span Programs. In Proceedings of the Structure in Complexity Theory Conference, pages 102–111, 1993. doi:10.1109/SCT.1993.336536.
\bibitem[50]{published-50} Ker-I Ko. On the Complexity of Learning Minimum Time-Bounded Turing Machines. SIAM J. Comput., 20(5):962–986, 1991. doi:10.1137/0220059.
\bibitem[51]{published-51} Andrei N Kolmogorov. Three approaches to the quantitative definition of information. Problems of information transmission, 1(1):1–7, 1965.
\bibitem[52]{published-52} Troy Lee and Andrei E. Romashchenko. Resource bounded symmetry of information revisited. Theor. Comput. Sci., 345(2-3):386–405, 2005. doi:10.1016/j.tcs.2005.07.017.
\bibitem[53]{published-53} Leonid A. Levin. The Tale of One-Way Functions. Probl. Inf. Transm., 39(1):92–103, 2003. doi:10.1023/A:1023634616182.
\bibitem[54]{published-54} Leonid A. Levin. How do we succeed in tasks like proving Fermat’s Theorem or predicting the Higgs boson?, 2021. An invited talk given at STOC 2021. The transcript is available at https://www.cs.bu.edu/fac/lnd/expo/stoc21/txt.pdf.
\bibitem[55]{published-55} Ming Li and Paul M. B. Vitányi. An Introduction to Kolmogorov Complexity and Its Applications, Third Edition. Texts in Computer Science. Springer, 2008. doi:10.1007/ 978-0-387-49820-1.
\bibitem[56]{published-56} Yanyi Liu and Rafael Pass. On One-way Functions from NP-Complete Problems. Electron. Colloquium Comput. Complex., 28:59, 2021. URL: https://eccc.weizmann.ac.il/report/ 2021/059.
\bibitem[57]{published-57} Luc Longpré and Sarah Mocas. Symmetry of Information and One-Way Functions. Inf. Process. Lett., 46(2):95–100, 1993. doi:10.1016/0020-0190(93)90204-M.
\bibitem[58]{published-58} Luc Longpré and Osamu Watanabe. On Symmetry of Information and Polynomial Time Invertibility. Inf. Comput., 121(1):14–22, 1995. doi:10.1006/inco.1995.1120.
\bibitem[59]{published-59} Zhenjian Lu and Igor Carboni Oliveira. An Efficient Coding Theorem via Probabilistic Representations and Its Applications. In Proceedings of the International Colloquium on Automata, Languages, and Programming (ICALP), pages 94:1–94:20, 2021. doi:10.4230/ LIPIcs.ICALP.2021.94.
\bibitem[60]{published-60} Daniele Micciancio and Oded Regev. Worst-Case to Average-Case Reductions Based on Gaus- sian Measures. SIAM J. Comput., 37(1):267–302, 2007. doi:10.1137/S0097539705447360.
\bibitem[61]{published-61} Noam Nisan and Avi Wigderson. Hardness vs Randomness. J. Comput. Syst. Sci., 49(2):149– 167, 1994. doi:10.1016/S0022-0000(05)80043-1.
\bibitem[62]{published-62} Igor Carboni Oliveira. Randomness and Intractability in Kolmogorov Complexity. In Proceed- ings of the International Colloquium on Automata, Languages, and Programming (ICALP), pages 32:1–32:14, 2019. doi:10.4230/LIPIcs.ICALP.2019.32.
\bibitem[63]{published-63} Igor Carboni Oliveira and Rahul Santhanam. Pseudodeterministic constructions in subex- ponential time. In Proceedings of the Symposium on Theory of Computing (STOC), pages 665–677, 2017. doi:10.1145/3055399.3055500.
\bibitem[64]{published-64} Ran Raz, Omer Reingold, and Salil P. Vadhan. Extracting all the Randomness and Reducing the Error in Trevisan’s Extractors. J. Comput. Syst. Sci., 65(1):97–128, 2002. doi:10.1006/ jcss.2002.1824.
\bibitem[65]{published-65} Hanlin Ren and Rahul Santhanam. Hardness of KT Characterizes Parallel Cryptography. In Proceedings of the Computational Complexity Conference (CCC), pages 35:1–35:58, 2021. doi:10.4230/LIPIcs.CCC.2021.35.
\bibitem[66]{published-66} Detlef Ronneburger. Kolmogorov Complexity and Derandomization. PhD thesis, Rutgers UniversityDept. of Computer Science Laboratory for Computer Sci. Research Hill Center, NJ, USA, 2004.
\bibitem[67]{published-67} Adi Shamir. How to Share a Secret. Commun. ACM, 22(11):612–613, 1979. doi:10.1145/ 359168.359176.
\bibitem[68]{published-68} Amnon Ta-Shma, Christopher Umans, and David Zuckerman. Lossless Condensers, Un- balanced Expanders, And Extractors. Combinatorica, 27(2):213–240, 2007. doi:10.1007/ s00493-007-0053-2.
\bibitem[69]{published-69} Boris A. Trakhtenbrot. A Survey of Russian Approaches to Perebor (Brute-Force Searches) Algorithms. IEEE Annals of the History of Computing, 6(4):384–400, 1984. doi:10.1109/ MAHC.1984.10036.
\bibitem[70]{published-70} Luca Trevisan. Extractors and pseudorandom generators. J. ACM, 48(4):860–879, 2001. doi:10.1145/502090.502099.
\bibitem[71]{published-71} Luca Trevisan and Salil P. Vadhan. Pseudorandomness and Average-Case Complexity Via Uniform Reductions. Computational Complexity, 16(4):331–364, 2007. doi:10.1007/ s00037-007-0233-x.
\bibitem[72]{published-72} Emanuele Viola. On Constructing Parallel Pseudorandom Generators from One-Way Functions. In Proceedings of the Conference on Computational Complexity (CCC), pages 183–197, 2005. doi:10.1109/CCC.2005.16.
\bibitem[73]{published-73} Emanuele Viola. The complexity of constructing pseudorandom generators from hard functions. Computational Complexity, 13(3-4):147–188, 2005. doi:10.1007/s00037-004-0187-1.
\bibitem[74]{published-74} AK Zvonkin and LA Levin. The complexity of finite objects and the algorithmic concepts of randomness and information. UMN (Russian Math. Surveys), 25(6):83–124, 1970.
\end{thebibliography}
\end{document}
